\documentclass[a4paper,12pt]{article}

\usepackage[utf8]{inputenc}
\usepackage[T2A]{fontenc}
\usepackage[russian]{babel}
\usepackage{amsmath}
\usepackage[a4paper,margin=2.5cm]{geometry}
\usepackage{enumitem}
\usepackage[hidelinks]{hyperref}

\begin{document}

\section{Задача}

Необходимо разработать программную систему для игры-симуляции инвестиций в строительство.
В игре несколько игроков управляют инвестиционным капиталом и вкладывают его в строительство жилых домов и супермаркетов. Игроки конкурируют между собой, стремясь получить максимальный общий капитал к концу игры.

\section{Параметры игры}

Количество игроков:
\[
2 \leq N \leq 5
\]

Количество игровых шагов:
\[
6 \leq M \leq 24
\]

Один игровой шаг соответствует одному месяцу.

В начале игры каждый игрок получает одинаковый начальный капитал. В исходном задании в качестве примера указан капитал в размере 37 млн у.е.

Основные изменяемые параметры игры:
\begin{itemize}
    \item количество игроков $N$;
    \item количество игровых месяцев $M$;
    \item стоимость строительства объектов;
    \item сроки строительства объектов;
    \item параметры объектов различных типов;
    \item параметры спроса и продаж.
\end{itemize}

\section{Объекты инвестирования}

Игрок может вкладывать средства в следующие объекты:
\begin{itemize}
    \item жилые дома;
    \item супермаркеты.
\end{itemize}

Для жилых домов могут использоваться различные типы строительства, например:
\begin{itemize}
    \item панельные дома;
    \item монолитные дома;
    \item кирпичные дома.
\end{itemize}

Для каждого типа объекта задаются стоимость строительства и срок строительства.

Например:
\begin{itemize}
    \item стоимость строительства дома --- от 8 млн;
    \item срок строительства дома --- от 7 месяцев;
    \item стоимость строительства супермаркета --- от 2,5 млн;
    \item срок строительства супермаркета --- от 5 месяцев.
\end{itemize}

\section{Ход игры}

В начале каждого игрового месяца для каждого игрока последовательно выполняются следующие действия:
\begin{enumerate}
    \item Выплачиваются текущие затраты на строительство объектов.
    \item Игрок получает доход от продаж жилья, выполненных в предыдущем месяце.
    \item Игрок получает прибыль от продаж товаров в действующих супермаркетах.
    \item Определяется объём жилья, выставляемого на продажу в текущем месяце.
    \item Определяется цена квадратного метра жилья.
    \item Определяются затраты на рекламу продаваемого жилья.
    \item Определяются затраты на рекламу товаров в построенных супермаркетах.
    \item Риэлторское агентство распределяет спрос на выставленное игроками жильё.
\end{enumerate}

\section{Строительство}

Стоимость строительства каждого объекта равномерно расходуется в течение всего срока его строительства.
Таким образом, строительство требует ежемесячного финансирования до момента завершения объекта.
Дополнительные супермаркеты могут повышать спрос на жильё, расположенное рядом.
Дополнительные жилые дома могут увеличивать уровень продаж в соседних супермаркетах.

\section{Продажа жилья}

Игрок определяет:
\begin{itemize}
    \item количество жилья, выставляемого на продажу;
    \item цену одного квадратного метра;
    \item затраты на рекламу.
\end{itemize}

Продажа жилья осуществляется через единое риэлторское агентство.

Риэлторское агентство учитывает:
\begin{itemize}
    \item текущий спрос;
    \item цену жилья;
    \item объём предложения;
    \item затраты на рекламу;
    \item другие параметры игровой модели.
\end{itemize}

Если общее предложение превышает спрос, в первую очередь удовлетворяется предложение жилья с более низкой ценой.

\section{Реклама жилья}

Затраты на рекламу влияют на объём продаж жилья.
Согласно исходному заданию, каждые 1 тыс., потраченные на рекламу, увеличивают объём продаваемого в текущем и следующем месяце жилья на 0,5\%.

\section{Реклама супермаркетов}

Игрок также может тратить деньги на рекламу товаров в уже построенных супермаркетах.
Согласно исходному заданию, каждые 0,5 тыс., потраченные на такую рекламу, увеличивают прибыль от продаж на 3\%.

\section{Спрос на жильё}

Спрос на жильё изменяется в зависимости от сезона года.

Обычный спрос:
\begin{itemize}
    \item увеличивается с весны;
    \item достигает максимального значения осенью.
\end{itemize}

Строительство дополнительных супермаркетов может дополнительно влиять на спрос на жильё.

\section{Продажи в супермаркетах}

Средний уровень продаж в супермаркетах также зависит от сезона.
Согласно исходному заданию, уровень продаж:
\begin{itemize}
    \item возрастает с осени;
    \item достигает максимального значения зимой.
\end{itemize}

На уровень продаж также может влиять количество жилых домов в микрорайоне.

\section{Общий капитал игрока}

К концу игры общий капитал игрока определяется с учётом:
\begin{itemize}
    \item стоимости строящихся объектов;
    \item стоимости непроданного жилья в построенных домах;
    \item стоимости построенных супермаркетов;
    \item денежных средств, полученных от продаж.
\end{itemize}

Стоимость непроданного жилья в построенных домах рассчитывается исходя из себестоимости одного квадратного метра.
Стоимость построенного супермаркета принимается равной потраченной на его строительство сумме, увеличенной на 60\%.

\section{Условия окончания игры}

Игра заканчивается после выполнения установленного количества игровых шагов $M$.
После окончания игры сравнивается общий капитал всех игроков.
Побеждает игрок, имеющий наибольший общий капитал.

\section{Требования к программной системе}

Программа должна включать:
\begin{itemize}
    \item модуль управления ходом моделирования;
    \item модуль контроля соблюдения правил игры;
    \item модуль работы риэлторского агентства;
    \item несколько модулей игроков, реализующих различные стратегии игры.
\end{itemize}

В процессе игры игрокам должна быть доступна информация:
\begin{itemize}
    \item о строящихся объектах других игроков;
    \item о построенных объектах других игроков;
    \item о проданном другими игроками жилье.
\end{itemize}

\end{document}
